These experiments establish the fixed-budget tradeoff and a conditional high-load goodput gain, but not a consistent improvement across full service metrics and loads. The frozen long-document confirmation does not beat its development-selected fixed policy. At moderate load, most actions collapse to the same cap; at higher load, the feedback target can be incompatible with the decode cost already present, so reducing prompt service can move delay into the waiting queue. The scope is one model, one GPU, in-process synchronous vLLM, fixed output lengths, and a small set of structured arrival traces. Model precision is unchanged, but different batch shapes can alter greedy outputs through floating-point execution; fixed output lengths hold service work constant and are not a semantic-equivalence claim. Host-observed delivery gaps omit network transport and do not isolate device-only latency. GPU clocks are not locked. A 100 ms gap SLO leaves many small-load requests qualified under every policy, so its goodput is often determined by full service duration. No hindsight choice of a fixed cap per trace segment is presented as an online algorithm. The high-load trace was used for fixed-budget tuning and is exploratory rather than an independent high-pressure holdout. The count-only comparator is a conservative simple rule, not a tuned count-based controller. Thus the value of elapsed time beyond a strong count rule remains unresolved; a MoE label supplies no novelty claim.
